
\section{\BookRecordIndex{T0144}格罗莫夫--威滕势}\label{sec:11-2}

上一节定义的\BookRecordIndex{T0228}量子杯积只用到了\BookRecordIndex{T0143}格罗莫夫--威滕不变量所提供的一小部分信息。本节定义一个
编码全部这些信息的形式幂级数，称为亏格零\BookRecordIndex{T0144}格罗莫夫--威滕势。我们将说明\BookRecordIndex{T0125}分裂公理如何化为
一组称为 WDVV 方程的三阶二次偏微分方程。本节的叙述受 Kontsevich 与 Manin
\srcite{216} 的工作启发。

设 $(M,\omega)$ 是紧\BookRecordIndex{T0058}半正\BookRecordIndex{T0358}辛流形，底环 $R$ 取为域
$\mathbb Q,\mathbb R,\mathbb C$ 之一。$R$ 的具体选择无关紧要，因为在\BookRecordIndex{T0058}半正情形，所有
至少有三个\BookRecordIndex{T0066}标记点的\BookRecordIndex{T0143}格罗莫夫--威滕不变量都是整数。与第~\ref{sec:11-1} 节一样，选取
$H^*(M)$ 的整数基 $e_0,\ldots,e_N$，使得 $e_0=1\in H^0(M)$，而每个基元素 $e_\nu$
次数纯齐；定义
\[
  g_{\nu\mu}:=\int_Me_\nu\smile e_\mu,
\]
并用 $g^{\nu\mu}$ 表示逆矩阵。接着引入 $N+1$ 个形式变量 $t_0,\ldots,t_N$，并定义以
$H^*(M)$ 为系数、关于 $t_0,\ldots,t_N$ 的线性多项式
\[
  a_t:=t_0e_0+t_1e_1+\cdots+t_Ne_N.
\]
$(M,\omega)$ 的\emph{\BookRecordIndex{T0144}格罗莫夫--威滕势}是变量 $t_0,\ldots,t_N$ 的形式幂级数，系数属于
Novikov 环 $\Lambda_\omega$；它定义为
\begin{equation}\label{eq:11-2-1}
  \Phi(t):=\sum_A\sum_p\frac1{p!}
  \operatorname{GW}_{A,p}^M(a_t,\ldots,a_t)e^A,
\end{equation}
其中形式求和遍历所有同调类 $A\in H_2(M)$ 和所有非负整数 $p$。其含义解释如下。

关键是把 $H^*(M)$ 看作带有偶坐标和奇坐标的\BookRecordIndex{T0075}超流形。若 $\deg(e_\nu)$ 为奇数，就称变量
$t_\nu$ 为奇变量，否则称为偶变量。在\BookRecordIndex{T0075}超流形中采用如下约定：奇变量反交换，偶变量交换。
因此
\begin{equation}\label{eq:11-2-2}
  t_\nu t_\mu=
  \begin{cases}
  -t_\mu t_\nu,&t_\nu,t_\mu\ \text{均为奇变量},\\
  t_\mu t_\nu,&\text{其他情形},
  \end{cases}
\end{equation}
并且
\begin{equation}\label{eq:11-2-3}
  t_\nu e_\mu=
  \begin{cases}
  -e_\mu t_\nu,&t_\nu,e_\mu\ \text{均为奇次数},\\
  e_\mu t_\nu,&\text{其他情形}.
  \end{cases}
\end{equation}
特别地，每个奇变量都满足 $(t_\nu)^2=0$。注意，$t_\nu$ 的交换关系恰好与
\BookRecordIndex{T0143}格罗莫夫--威滕不变量之输入 $e_\nu$ 的交换关系相同。此外，以任意交换环为系数、关于
$t_0,\ldots,t_N$ 的形式幂级数，都可写成
$\sum_\alpha\lambda_\alpha t^\alpha$，其中求和遍历整数向量
$\alpha=(\alpha_0,\ldots,\alpha_N)$，满足 $\alpha_\nu\ge0$，且只要 $t_\nu$ 为奇变量
就有 $\alpha_\nu\in\{0,1\}$。

为理解公式~\eqref{eq:11-2-1}，更仔细地考察其右端的项
$\operatorname{GW}_{A,p}^M(a_t,\ldots,a_t)$。它可写成
\begin{equation}\label{eq:11-2-4}
\begin{aligned}
  \frac1{p!}\operatorname{GW}_{A,p}^M(a_t,\ldots,a_t)
  &=
  \frac1{p!}\sum_{\nu_1,\ldots,\nu_p}
  \operatorname{GW}_{A,p}^M
  (t_{\nu_1}e_{\nu_1},\ldots,t_{\nu_p}e_{\nu_p})\\
  &=
  \sum_{\nu_1,\ldots,\nu_p}\frac{\varepsilon(\nu)}{p!}
  \operatorname{GW}_{A,p}^M(e_{\nu_1},\ldots,e_{\nu_p})
  t_{\nu_1}\cdots t_{\nu_p}\\
  &=
  \sum_{|\alpha|=p}\frac{\varepsilon(\alpha)}{\alpha!}
  \operatorname{GW}_{A,|\alpha|}^M(e^\alpha)t^\alpha.
\end{aligned}
\end{equation}
这里第一个和遍历所有满足 $0\le\nu_i\le N$ 的整数 $p$-元组
$(\nu_1,\ldots,\nu_p)$。符号
$\varepsilon(\nu)=\varepsilon(\alpha)\in\{\pm1\}$ 来自重新排列乘积
\[
  t_{\nu_1}e_{\nu_1}\cdots t_{\nu_p}e_{\nu_p}
  =\varepsilon(\nu)e_{\nu_1}\cdots e_{\nu_p}
  t_{\nu_1}\cdots t_{\nu_p}.
\]
因此
\[
  \varepsilon(\nu)=(-1)^{r(r+1)/2},
\]
其中 $r$ 是 $t_{\nu_1},\ldots,t_{\nu_p}$ 中奇变量的数目。
\sourcefootnote{原书此处把符号误印为
$\varepsilon(\nu)=r(r+1)/2$，并在习题 11.2.3 中再次把
$\varepsilon(\alpha)$ 误印为 $r_\alpha(r_\alpha+1)/2$，两处都漏掉了
$(-1)$ 及幂号；前文已明确 $\varepsilon\in\{\pm1\}$，且紧接着的定义唯一确定了补正。}
类似地，定义
\[
  \varepsilon(\alpha)=(-1)^{r(r+1)/2},\qquad
  r:=r_\alpha:=\#\{\nu\mid\alpha_\nu=1,\ t_\nu\ \text{为奇变量}\}.
\]
式~\eqref{eq:11-2-2} 中的符号约定是一种组合装置，它保证置换 $e_{\nu_i}$ 所引起的符号
变化与对 $t_{\nu_i}$ 作同一置换所引起的变化一致。于是得到~\eqref{eq:11-2-4} 的最后
一个表达式，其中
\[
  t^\alpha:=(t_0)^{\alpha_0}\cdots(t_N)^{\alpha_N},
  \qquad \alpha!:=\alpha_0!\cdots\alpha_N!,
\]
并且
\[
  \operatorname{GW}_{A,p}^M(e^\alpha)
  :=\operatorname{GW}_{A,p}^M(
  \underbrace{e_0,\ldots,e_0}_{\alpha_0},
  \underbrace{e_1,\ldots,e_1}_{\alpha_1},
  \ldots,
  \underbrace{e_N,\ldots,e_N}_{\alpha_N}).
\]
换言之，每个 $e_\nu$ 出现 $\alpha_\nu$ 次。\eqref{eq:11-2-4} 中最后一个和遍历所有满足
$|\alpha|:=\alpha_0+\cdots+\alpha_N=p$ 的非负整数向量
$\alpha=(\alpha_0,\ldots,\alpha_N)$。把~\eqref{eq:11-2-4} 代入~\eqref{eq:11-2-1}，
得到
\begin{equation}\label{eq:11-2-5}
  \Phi(t)=\sum_\alpha\sum_A\frac{\varepsilon(\alpha)}{\alpha!}
  \operatorname{GW}_{A,|\alpha|}^M(e^\alpha)e^At^\alpha.
\end{equation}
这里 $|\alpha|<3$ 的不变量按注~\ref{rem:7-5-2} 理解。维数条件
\begin{equation}\label{eq:11-2-6}
  \sum_{\nu=0}^N\alpha_\nu\deg(e_\nu)
  =\dim M+2c_1(A)+2|\alpha|-6
\end{equation}
保证级数~\eqref{eq:11-2-5} 中 $t^\alpha$ 的系数是 $\Lambda_\omega$ 中次数纯齐的元素，
其次数为
\[
  d_\alpha:=\sum_{\nu=0}^N\alpha_\nu\bigl(\deg(e_\nu)-2\bigr)+6-\dim M.
\]
换言之，若赋予每个 $t_\nu$ 次数 $2-\deg(e_\nu)$，则幂级数~\eqref{eq:11-2-5} 是次数
$6-\dim M$ 的齐次级数。

\begin{example}\label{ex:11-2-1}
把复系数上同调类 $a_t\in H^*(\mathbb{CP}^n)$ 写成形式线性组合
\[
  a_t:=t_0+t_1c+t_2c^2+\cdots+t_nc^n,
\]
其中 $t:=(t_0,\ldots,t_n)$，并用 $q$ 表示与
$H_2(\mathbb{CP}^n)$ 的生成元 $L$ 对应的变量。第~\ref{sec:7-5} 节已经看到，
$\mathbb{CP}^1$ 的\BookRecordIndex{T0144}格罗莫夫--威滕势为
\[
  \Phi(t_0,t_1)=\frac12t_0^2t_1+qe^{t_1},
\]
而 $\mathbb{CP}^2$ 的\BookRecordIndex{T0144}格罗莫夫--威滕势具有形式
\[
  \Phi^{\mathbb{CP}^2}(t_0,t_1,t_2)
  =\frac12(t_0^2t_2+t_0t_1^2)
  +\sum_{d>0}N_d\frac{(t_2)^{3d-1}}{(3d-1)!}q^de^{dt_1},
\]
其中 $N_d$ 是通过 $3d-1$ 个通有点的次数 $d$、亏格零曲线数。
\end{example}

下面适合这样选基：$e_1,\ldots,e_m$ 构成 $H^2(M)$ 的整数基。用
$A_1,\ldots,A_m$ 表示 $H_2(M)$ 的对偶基，使得
\[
  \int_{A_j}e_i=\delta_{ij},\qquad i,j=1,\ldots,m.
\]
于是，可通过同构 $\mathbb Z^m\to H_2(M):d\mapsto A_d$ 把 $H_2(M)$ 与 $\mathbb Z^m$
等同，其中
\[
  A_d:=d_1A_1+\cdots+d_mA_m,\qquad
  d=(d_1,\ldots,d_m)\in\mathbb Z^m.
\]
还引入形式变量 $q_i:=e^{A_i}$。于是，\BookRecordIndex{T0144}格罗莫夫--威滕势可写成 $q_i,t_\nu$ 的形式幂级数：
\begin{equation}\label{eq:11-2-7}
  \Phi(q,t):=\sum_\alpha\sum_d\frac{\varepsilon(\alpha)}{\alpha!}
  \operatorname{GW}_{A_d,|\alpha|}^M(e^\alpha)q^dt^\alpha.
\end{equation}
式~\eqref{eq:11-2-2} 保证 $\Phi$ 关于每个奇变量 $t_\nu$ 都是仿射的。

\begin{remark}\label{rem:11-2-2}
对奇变量求幂级数的导数时必须小心。例如，若 $t_\nu,t_\mu$ 是两个不同奇变量，则乘积
$t_\nu t_\mu$ 与 $t_\mu t_\nu$ 对 $t_\nu$ 的偏导数应符号相反，且分别等于
$\pm t_\mu$。我们采用如下约定：单项式 $t^\alpha$ 对奇变量 $t_\nu$ 的偏导数，通过先把
$t_\nu$ 移到最前面再删去它来定义。因此
\[
  \frac{\partial}{\partial t_\nu}t^\alpha
  :=\varepsilon(\nu,\alpha)\alpha_\nu\,t^\alpha/t_\nu
\]
只要 $\alpha_\nu>0$。\sourcefootnote{原书把此式的因子印作
$(-1)^{\varepsilon(\nu,\alpha)}$，而紧接着的~\eqref{eq:11-2-8} 又把
$\varepsilon(\nu,\alpha)$ 定义为取值于 $\{\pm1\}$ 的符号，两式不\BookRecordIndex{T0348}相容。也可把
$\varepsilon$ 定义成模二奇偶数并保留外\BookRecordIndex{T0072}层 $(-1)$；本译稿依本节前后一贯的“符号”约定，
直接乘以 $\varepsilon(\nu,\alpha)$。}
这里
\[
  t^\alpha/t_\nu:=t^{\alpha'},\qquad
  \alpha'_\mu:=
  \begin{cases}
  \alpha_\mu-1,&\mu=\nu,\\
  \alpha_\mu,&\mu\ne\nu,
  \end{cases}
\]
并且
\begin{equation}\label{eq:11-2-8}
  \varepsilon(\nu,\alpha):=
  (-1)^{\deg(e_\nu)\cdot
  \#\{\mu<\nu\mid t_\mu\ \text{为奇变量},\ \alpha_\mu=1\}}.
\end{equation}
在此约定下，对两个不同奇变量的偏导数反交换；其他所有情形的偏导数都交换。
\end{remark}

\begin{exercise}\label{ex:11-2-3}
给定整数向量 $\alpha$，用 $r_\alpha$ 表示满足 $\alpha_\nu=1$ 的奇变量数，令
\[
  \varepsilon(\alpha):=(-1)^{r_\alpha(r_\alpha+1)/2},
\]
并设 $\alpha'$ 是把 $\alpha_\nu$ 减去一所得的向量。
\begin{enumerate}[label=\textnormal{(\roman*)}]
\item 证明
\[
  \varepsilon(\alpha')
  =(-1)^{r_\alpha|e_\nu|}\varepsilon(\alpha),
\]
其中 $|e_\nu|$ 表示 $e_\nu$ 的次数。因此
\[
  \varepsilon(\alpha)t_\nu t^{\alpha'}
  =(-1)^{|e_\nu|}\varepsilon(\alpha')t^{\alpha'}t_\nu.
\]
\item 假定 $\alpha,\beta$ 没有共同的奇变量。证明
\[
  \varepsilon(\alpha)\varepsilon(\beta)
  =(-1)^{r_\alpha r_\beta}\varepsilon(\alpha+\beta).
\]
\end{enumerate}
\end{exercise}

下一个定理说明\BookRecordIndex{T0143}格罗莫夫--威滕不变量的公理如何化为\BookRecordIndex{T0144}格罗莫夫--威滕势的偏微分方程。
（基本类）公理化为\BookRecordIndex{T0345}弦方程，（除子）公理化为\BookRecordIndex{T0077}除子方程，（分裂）公理化为 WDVV 方程
（Witten--Dijkgraaf--Verlinde--Verlinde）。在当前背景下，这些方程非常简单；在高亏格
情形则有趣得多，见 Manin \srcite{286}。

\begin{theorem}\label{thm:11-2-4}
设 $(M,\omega)$ 是闭\BookRecordIndex{T0058}半正\BookRecordIndex{T0358}辛流形。则\BookRecordIndex{T0144}格罗莫夫--威滕势 $\Phi=\Phi^M$ 满足\BookRecordIndex{T0345}弦方程
\begin{equation}\label{eq:11-2-9}
  \partial_{t_0}\Phi(q,t)=\frac12\int_Ma_t\smile a_t,
\end{equation}
\BookRecordIndex{T0077}除子方程
\begin{equation}\label{eq:11-2-10}
  \partial_{t_i}\Phi(q,t)
  =q_i\partial_{q_i}\Phi(q,t)
  +\frac12\int_Me_i\smile a_t\smile a_t,
  \qquad i=1,\ldots,m,
\end{equation}
以及 WDVV 方程
\begin{equation}\label{eq:11-2-11}
\begin{aligned}
  \sum_{\nu,\mu}
  \partial_{t_i}\partial_{t_j}\partial_{t_\nu}\Phi\,
  g^{\nu\mu}\,
  \partial_{t_\mu}\partial_{t_k}\partial_{t_\ell}\Phi
  &=
  \varepsilon_{ijk}\sum_{\nu,\mu}
  \partial_{t_j}\partial_{t_k}\partial_{t_\nu}\Phi\,
  g^{\nu\mu}\,
  \partial_{t_\mu}\partial_{t_i}\partial_{t_\ell}\Phi,
\end{aligned}
\end{equation}
其中
\[
  \varepsilon_{ijk}:=
  (-1)^{\deg(e_i)(\deg(e_j)+\deg(e_k))}.
\]
\end{theorem}

\begin{proof}
首先写
\[
  \Phi(q,t)=\sum_d\Phi_d(t)q^d,
\]
其中 $\Phi_d$ 是变量 $t_0,\ldots,t_N$ 中具有有理系数的形式幂级数：
\[
  \Phi_d(t):=\sum_p\frac1{p!}
  \operatorname{GW}_{A_d,p}^M(a_t,\ldots,a_t)
  =\sum_\alpha\frac{\varepsilon(\alpha)}{\alpha!}
  \operatorname{GW}_{A_d,|\alpha|}^M(e^\alpha)t^\alpha.
\]
采用注~\ref{rem:11-2-2} 的符号约定和记号，$\Phi_d$ 对 $t_\nu$ 的偏导数为
\begin{equation}\label{eq:11-2-12}
\begin{aligned}
  \partial_{t_\nu}\Phi_d(t)
  &=
  \sum_{\alpha_\nu\ne0}
  \frac{\varepsilon(\alpha)\alpha_\nu}{\alpha!}
  \operatorname{GW}_{A_d,|\alpha|}^M(e_\nu,e^{\alpha'})t^{\alpha'}\\
  &=
  \sum_\alpha(-1)^{|e_\nu|}
  \frac{\varepsilon(\alpha)}{\alpha!}
  \operatorname{GW}_{A_d,|\alpha|+1}^M(e^\alpha,e_\nu)t^\alpha.
\end{aligned}
\end{equation}
第一个等式使用如下事实：对 $e_\nu,t_\nu$ 作同一置换不会改变符号；第二个等式使用
习题~\ref{ex:11-2-3}(i)。

先证明\BookRecordIndex{T0345}弦方程。在~\eqref{eq:11-2-12} 中取 $\nu=0$，得到
\[
\begin{aligned}
  \partial_{t_0}\Phi(q,t)
  &=
  \sum_\alpha\sum_d\frac{\varepsilon(\alpha)}{\alpha!}
  \operatorname{GW}_{A_d,|\alpha|+1}^M(e^\alpha,1)q^dt^\alpha\\
  &=
  \sum_{|\alpha|=2}\frac{\varepsilon(\alpha)}{\alpha!}
  \operatorname{GW}_{0,3}^M(1,e^\alpha)t^\alpha\\
  &=
  \frac12\sum_{i,j}\int_Mt_ie_i\smile t_je_j\\
  &=
  \frac12\int_Ma_t\smile a_t.
\end{aligned}
\]
第二个等式由（基本类）和（零类）公理\BookRecordIndex{T0326}推出。

再证明\BookRecordIndex{T0077}除子方程。在~\eqref{eq:11-2-12} 中取 $\nu=i\in\{1,\ldots,m\}$，得到
\[
\begin{aligned}
  \partial_{t_i}\Phi(q,t)
  &=
  \sum_\alpha\sum_d\frac{\varepsilon(\alpha)}{\alpha!}
  \operatorname{GW}_{A_d,|\alpha|+1}^M(e^\alpha,e_i)q^dt^\alpha\\
  &=
  \sum_\alpha\sum_{d\ne0}\frac{\varepsilon(\alpha)}{\alpha!}
  \operatorname{GW}_{A_d,|\alpha|}^M(e^\alpha)
  \left(\int_{A_d}e_i\right)q^dt^\alpha\\
  &\quad+
  \sum_{|\alpha|=2}\frac{\varepsilon(\alpha)}{\alpha!}
  \operatorname{GW}_{0,3}^M(e_i,e^\alpha)t^\alpha\\
  &=
  q_i\partial_{q_i}\Phi(q,t)
  +\frac12\int_Me_i\smile a_t\smile a_t.
\end{aligned}
\]
第二个等式来自（除子）和（零类）公理，第三个等式则使用
$\int_{A_d}e_i=d_i$ 以及对每个 $i$ 都有 $d_iq^d=q_i\partial_{q_i}q^d$。

最后证明 WDVV 方程。迭代恒等式~\eqref{eq:11-2-12}，得到
\[
  \partial_{t_i}\partial_{t_j}\partial_{t_\nu}\Phi_d(t)
  =
  \sum_\alpha(-1)^{r_\alpha}
  \frac{\varepsilon(\alpha)}{\alpha!}
  \operatorname{GW}_{A_d,|\alpha|+3}^M
  (e^\alpha,e_i,e_j,e_\nu)t^\alpha.
\]
这里用到了如下事实：\BookRecordIndex{T0143}格罗莫夫--威滕不变量只有在所代入各类的总次数为偶数时才可能非零，
故 $r_\alpha$ 与 $|e_i|+|e_j|+|e_\nu|$ 奇偶性相同。用
\[
  \Phi_{ijk\ell}(t):=\sum_d\Phi_{ijk\ell,d}(t)q^d
\]
表示~\eqref{eq:11-2-11} 左端。由刚才的观察，
\[
\begin{aligned}
  \Phi_{ijk\ell,d}(t)
  &=
  \sum_{d'+d''=d}\sum_{\nu,\mu}\sum_{\alpha,\beta}
  (-1)^{r_\alpha+r_\beta}
  \frac{\varepsilon(\alpha)\varepsilon(\beta)}{\alpha!\beta!}\\
  &\qquad\qquad{}\times
  \operatorname{GW}_{A_{d'}}^M(e^\alpha,e_i,e_j,e_\nu)g^{\nu\mu}
  \operatorname{GW}_{A_{d''}}^M(e^\beta,e_\mu,e_k,e_\ell)t^\alpha t^\beta.
\end{aligned}
\]
这里略去了记号中的\BookRecordIndex{T0066}标记点数，因为上下文已明确：两个不变量分别有 $|\alpha|+3$ 和
$|\beta|+3$ 个\BookRecordIndex{T0066}标记点。现在用恒等式
\[
  t^\alpha t^\beta=\varepsilon(\alpha,\beta)t^{\alpha+\beta}
\]
定义 $\varepsilon(\alpha,\beta)\in\{\pm1\}$，并回忆习题~\ref{ex:11-2-3}(ii) 给出的
\[
  \varepsilon(\alpha+\beta)
  =(-1)^{r_\alpha r_\beta}\varepsilon(\alpha)\varepsilon(\beta).
\]
于是
\[
  \Phi_{ijk\ell,d}(t)=\sum_\gamma\frac1{\gamma!}
  \Phi_{ijk\ell,d,\gamma}t^\gamma,
\]
其中
\[
\begin{aligned}
  \Phi_{ijk\ell,d,\gamma}
  &:=
  \sum_{\substack{d'+d''=d\\ \alpha+\beta=\gamma}}
  \sum_{\nu,\mu}
  (-1)^{r_\alpha+r_\beta+r_\alpha r_\beta}
  \varepsilon(\gamma)\varepsilon(\alpha,\beta)
  \frac{\gamma!}{\alpha!\beta!}\\
  &\qquad{}\times
  \operatorname{GW}_{A_{d'}}^M(e^\alpha,e_i,e_j,e_\nu)g^{\nu\mu}
  \operatorname{GW}_{A_{d''}}^M(e^\beta,e_\mu,e_k,e_\ell)\\
  &=
  \sum_{\substack{d'+d''=d\\ \alpha+\beta=\gamma}}
  \sum_{\nu,\mu}
  (-1)^{r_\alpha+r_\beta+r_\alpha r_\beta}
  \varepsilon(\gamma)\varepsilon(\alpha,\beta)
  \frac{\gamma!}{\alpha!\beta!}\\
  &\qquad{}\times
  \operatorname{GW}_{A_{d'}}^M(e^\alpha,e_i,e_j,e_\nu)g^{\nu\mu}
  \operatorname{GW}_{A_{d''}}^M(e_\mu,e_k,e_\ell,e^\beta)\\
  &=
  \sum_{\substack{d'+d''=d\\ \alpha+\beta=\gamma}}
  \sum_{\nu,\mu}
  (-1)^{r_\alpha r_\gamma}
  \varepsilon(\gamma)\varepsilon(\alpha,\beta)
  \frac{\gamma!}{\alpha!\beta!}\\
  &\qquad{}\times
  \operatorname{GW}_{A_{d'}}^M(e^\alpha,e_i,e_j,e_\nu)g^{\nu\mu}
  \operatorname{GW}_{A_{d''}}^M(e_\mu,e_k,e_\ell,e^\beta)\\
  &=
  \varepsilon(\gamma)
  \operatorname{GW}_{A_d,|\gamma|+4}^{M,\{1,2,3,4\}}
  (e_i,e_j,e_k,e_\ell,e^\gamma).
\end{aligned}
\]
最后一个等式来自定理~\ref{thm:7-5-10} 对四个\BookRecordIndex{T0150}固定标记点给出的（分裂）公理。为理解符号，
注意 $e^\alpha$ 越过 $e_i,e_j,e_\nu,e_\mu,e_k,e_\ell$ 六项时，其总奇偶性与
$e^\gamma$ 相同，这抵消符号 $(-1)^{r_\alpha r_\gamma}$。把 $e^\alpha,e^\beta$ 这对输入
合并为 $e^\gamma$ 时，又抵消符号 $\varepsilon(\alpha,\beta)$。明白这一点后，
方程~\eqref{eq:11-2-11} 就由\BookRecordIndex{T0143}格罗莫夫--威滕不变量对前三个变量的分次对称性\BookRecordIndex{T0326}推出。定理得证。
\end{proof}

\begin{exercise}\label{ex:11-2-5}
证明\BookRecordIndex{T0345}弦方程可写成
\begin{equation}\label{eq:11-2-13}
  \partial_{t_0}\partial_{t_i}\partial_{t_j}\Phi=g_{ij}.
\end{equation}
\emph{提示：}变量 $t_i,t_j$ 要么都是偶变量，要么都是奇变量。
\end{exercise}

\begin{exercise}\label{ex:11-2-6}
令 $\mathcal E$ 为形式微分算子
\[
  \mathcal E:=
  \sum_{i=1}^mc^iq_i\partial_{q_i}
  +\sum_{\nu=0}^N
  \left(1-\frac{\deg(e_\nu)}2\right)t_\nu\partial_{t_\nu},
\]
其中第一陈类 $c_1(TM)=\sum_i c^ie_i$。证明
\[
  \mathcal E\Phi=(3-n)\Phi,\qquad 2n:=\dim M.
\]
这是维数条件~\eqref{eq:11-2-6} 的结果。
\end{exercise}

第~\ref{sec:7-5} 节详细讨论了 $\mathbb{CP}^n$ 的\BookRecordIndex{T0144}格罗莫夫--威滕势
$\Phi^{\mathbb{CP}^n}$。特别地，我们在命题~\ref{prop:7-5-11} 中证明，
$\Phi^{\mathbb{CP}^2}$ 由 WDVV 方程和初始条件
\[
  \operatorname{GW}_{L,3}^{\mathbb{CP}^2}(c^2,c^2,c^2)=1
\]
确定。虽然这个结果可以推广（见下面定理 11.5.8），但它并非总成立；比较第 11.3.4 小节
关于 Calabi--Yau 流形的讨论。第 11.5 节将更详细地探讨 WDVV 方程的几何意义，并说明它
如何在\BookRecordIndex{T0358}辛流形的上同调上导出 \BookRecordIndex{T0015}Dubrovin 联络等有趣结构。该节将只考察偶数次上同调，从而
避免颇为棘手的符号问题干扰讨论。
